\documentclass[11pt,a4paper,reqno]{amsart}

\usepackage[T1]{fontenc}
\usepackage[utf8]{inputenc}
\usepackage{lmodern}
\usepackage{amsmath,amssymb,mathtools}
\usepackage{microtype}
\usepackage[colorlinks=true,linkcolor=blue,citecolor=blue,urlcolor=blue]{hyperref}
\hypersetup{
 pdftitle={Unfolding the arcsine in generalized Hilbert determinants},
 pdfauthor={The authors}
}

\newtheorem{theorem}{Theorem}[section]
\newtheorem{proposition}[theorem]{Proposition}
\newtheorem{lemma}[theorem]{Lemma}
\newtheorem{corollary}[theorem]{Corollary}
\theoremstyle{remark}
\newtheorem{remark}[theorem]{Remark}

\newcommand{\Nzero}{\mathbb N_0}
\newcommand{\C}{\mathbb C}
\newcommand{\R}{\mathbb R}
\newcommand{\D}{\mathbb D}
\newcommand{\T}{\mathbb T}
\newcommand{\Hh}{\mathcal H}
\newcommand{\Tr}{\operatorname{Tr}}
\newcommand{\Ker}{\operatorname{Ker}}
\newcommand{\Ran}{\operatorname{Ran}}
\newcommand{\dist}{\operatorname{dist}}
\newcommand{\id}{I}
\newcommand{\cS}{\mathfrak S_1}

\title[Unfolding the arcsine]{Unfolding the arcsine in generalized Hilbert determinants}
\author{The authors}
\date{}

\begin{document}

\begin{abstract}
For $N\in\mathbb N$ and
$\delta\in\mathbb R\setminus\{-1,-2,\ldots\}$, let
\[
 H_N^\delta
 =\left(
 \frac{\sin((1+\delta)\pi)}{\pi(j+k+1+\delta)}
 \right)_{j,k=0}^{N-1},
 \qquad
 D_N^\pm(\delta):=\det(I_N\pm H_N^\delta).
\]
Write
\[
 q_\delta:=\frac1\pi\arcsin\bigl(\sin(-\pi\delta)\bigr)
 \in\left[-\frac12,\frac12\right]
\]
for the principal inverse-sine phase, and define the effective phase
\[
 \theta_\delta:=
 \begin{cases}
  -\delta,&\delta\le\frac12,\\
  q_\delta,&\delta\ge\frac12.
 \end{cases}
\]
The two definitions agree at $\delta=\frac12$.  If
$\delta\notin\frac12+\mathbb Z$, then
\[
 D_N^\pm(\delta)\asymp_\delta
 N^{-(\theta_\delta^2\mp\theta_\delta)/2}.
\]
At the half-integers the same powers hold up to subpolynomial factors:
\[
 D_N^\pm(\delta)
 =N^{-(\theta_\delta^2\mp\theta_\delta)/2+o(1)}.
\]
In particular, for $\delta\le\frac12$ away from the negative integers and
half-integers,
\[
 \det\bigl(I_N-(H_N^\delta)^2\bigr)\asymp_\delta N^{-\delta^2},
\]
while at the half-integers the right-hand side is multiplied by $N^{o(1)}$.

The proof separates the principal-branch Hilbert determinant from the
contribution of the finitely many eigenvectors of the infinite normalized
generalized Hilbert matrix at $+1$ and $-1$.  A finite-section lemma compares
the resulting small determinant, in Loewner order, with the Gram determinant
of the omitted eigenvector tails.  Their number and parity recover exactly
the branch index lost when $-\pi\delta$ is replaced by
$\arcsin(\sin(-\pi\delta))$.  At the critical half-integers the continuous
spectral gap closes; a separate Schur-complement argument retains the same
power of $N$, but presently only with an $N^{o(1)}$ remainder.
\end{abstract}

\maketitle

\section{Introduction and statement of the result}

Let $\mathbb N_0:=\{0,1,2,\ldots\}$.  For $N\in\mathbb N$ and
$\delta\in\R$, define the finite generalized Hilbert matrix
\begin{equation}\label{eq:def-Hdelta}
 H_N^\delta
 :=\left(
 \frac{\sin((1+\delta)\pi)}{\pi(j+k+1+\delta)}
 \right)_{j,k=0}^{N-1}.
\end{equation}
The literal formula is undefined at $\delta=-1,-2,\ldots$, because one or more
entries have the form $0/0$.  These parameters are excluded until
Section~\ref{sec:exceptional}, where the entrywise analytic continuation is
described.

The determinants of interest are
\begin{equation}\label{eq:def-Dpm}
 D_N^\pm(\delta):=\det(\id_N\pm H_N^\delta).
\end{equation}
Paired signs are read vertically throughout the paper.  Thus $D_N^+$ denotes
$\det(\id_N+H_N^\delta)$ and $D_N^-$ denotes
$\det(\id_N-H_N^\delta)$.

These determinants form a finite-section model for power-law determinant
asymptotics of the type that occurs in Anderson's orthogonality catastrophe.
Our aim is to determine their powers of $N$ up to bounded multiplicative
factors whenever the parameter is noncritical, and up to subpolynomial factors
at the critical half-integers.  For positive sequences $a_N,b_N$, the notation
\[
 a_N\asymp_\delta b_N
\]
means that $c_\delta b_N\le a_N\le C_\delta b_N$ for some constants
$0<c_\delta\le C_\delta<\infty$ independent of $N$.  Equivalently,
$\log a_N=\log b_N+O_\delta(1)$.  We write
$a_N=b_NN^{o(1)}$ when
\[
 \log(a_N/b_N)=o(\log N).
\]
This permits subpolynomial factors, such as fixed powers of $\log N$.

Set
\begin{equation}\label{eq:def-sq}
 s_\delta:=\sin((1+\delta)\pi)=\sin(-\pi\delta),
 \qquad
 q_\delta:=\frac1\pi\arcsin(s_\delta)
 \in\left[-\frac12,\frac12\right],
\end{equation}
where $\arcsin:[-1,1]\to[-\pi/2,\pi/2]$ is the principal inverse sine.  Define
also
\begin{equation}\label{eq:effective-phase}
 \theta_\delta:=
 \begin{cases}
  -\delta,&\delta\le\frac12,\\[1mm]
  q_\delta,&\delta\ge\frac12.
 \end{cases}
\end{equation}
The definitions agree at $\delta=\frac12$, because $q_{1/2}=-1/2$.
The standard Hilbert determinant naturally produces the folded phase
$q_\delta$.  The main point of the paper is that for $\delta< -1/2$ the edge
eigenvectors of the infinite operator unfold this phase to $-\delta$.  No such
unfolding occurs for larger positive $\delta$.

The parameter $\delta$ is visible without any inverse trigonometric function in
a simple Gram representation.  Put
\[
 a:=\frac{1+\delta}{2},
 \qquad
 E_N:=\operatorname{diag}(1,-1,1,-1,\ldots),
\]
and, for the standard basis $(e_j)_{j=0}^{N-1}$ of $\C^N$, define
$C_{N,a},S_{N,a}:\C^N\to L^2(0,\pi)$ by
\[
 C_{N,a}e_j=\sqrt{\frac2\pi}\cos((j+a)x),
 \qquad
 S_{N,a}e_j=\sqrt{\frac2\pi}\sin((j+a)x).
\]
The product-to-sum identities give
\begin{equation}\label{eq:gram-representation}
 C_{N,a}^*C_{N,a}=E_N(\id_N+H_N^\delta)E_N,
 \qquad
 S_{N,a}^*S_{N,a}=E_N(\id_N-H_N^\delta)E_N.
\end{equation}
Indeed,
\[
 \frac2\pi\int_0^\pi \cos((j+a)x)\cos((k+a)x)\,dx
 =\delta_{jk}+(-1)^{j+k}
 \frac{\sin((1+\delta)\pi)}{\pi(j+k+1+\delta)},
\]
and the corresponding sine integral has the opposite sign in the second term.
If $\delta$ is not a negative integer, the frequencies are distinct up to sign
and none of the sine functions vanishes identically.  The two trigonometric
systems are therefore linearly independent.  Consequently,
\begin{equation}\label{eq:positive-determinants}
 D_N^\pm(\delta)>0
 \qquad
 (\delta\notin\{-1,-2,\ldots\}),
\end{equation}
so all logarithms below are real.  Formula~\eqref{eq:gram-representation} also
shows that the unfurled parameter is present from the outset as the frequency
shift $a=(1+\delta)/2$.

Define the pole and critical sets
\begin{equation}\label{eq:exceptional-sets}
 \mathcal P:=\{-1,-2,\ldots\},
 \qquad
 \mathcal C:=\frac12+\mathbb Z.
\end{equation}

\begin{theorem}[Main result]\label{thm:main}
Let $\delta\in\mathbb R\setminus\mathcal P$.

\begin{enumerate}
\item If $\delta\notin\mathcal C$, then
\begin{equation}\label{eq:strong-main}
 D_N^\pm(\delta)\asymp_\delta
 N^{-(\theta_\delta^2\mp\theta_\delta)/2}.
\end{equation}
Equivalently,
\begin{equation}\label{eq:strong-main-log}
 \log D_N^\pm(\delta)
 =-\frac{\theta_\delta^2\mp\theta_\delta}{2}\log N+O_\delta(1).
\end{equation}

\item If $\delta\in\mathcal C$, then
\begin{equation}\label{eq:critical-main}
 D_N^\pm(\delta)
 =N^{-(\theta_\delta^2\mp\theta_\delta)/2+o(1)}.
\end{equation}
\end{enumerate}
Consequently,
\begin{align}
 \det\bigl(I_N-(H_N^\delta)^2\bigr)
 &\asymp_\delta N^{-\theta_\delta^2},
 &&\delta\notin\mathcal C,
 \label{eq:product-global-strong}\\
 \det\bigl(I_N-(H_N^\delta)^2\bigr)
 &=N^{-\theta_\delta^2+o(1)},
 &&\delta\in\mathcal C.
 \label{eq:product-global-critical}
\end{align}
\end{theorem}

\begin{remark}[The critical sign]
At a half-integer, only one of the two determinants requires the
subpolynomial remainder in \eqref{eq:critical-main}.  It is $D_N^-$ when
$s_\delta=1$ and $D_N^+$ when $s_\delta=-1$, namely the determinant for which
the continuous spectrum reaches zero.  For the opposite sign, the proof gives
an $O_\delta(1)$ remainder in the logarithm.
\end{remark}

The result takes a particularly simple form in the unfolded range.

\begin{corollary}[Recovery of the original phase]\label{cor:unfolded}
Let $\delta\le\frac12$ and $\delta\notin\mathcal P$.  If
$\delta\notin\mathcal C$, then
\begin{align}
 \det(I_N-H_N^\delta)
 &\asymp_\delta N^{-\delta(\delta-1)/2},
 \label{eq:minus-result}\\
 \det(I_N+H_N^\delta)
 &\asymp_\delta N^{-\delta(\delta+1)/2},
 \label{eq:plus-result}\\
 \det\bigl(I_N-(H_N^\delta)^2\bigr)
 &\asymp_\delta N^{-\delta^2}.
 \label{eq:square-result}
\end{align}
If $\delta\in\mathcal C$, the three relations remain true after replacing
$\asymp_\delta$ by equality up to a factor $N^{o(1)}$.
\end{corollary}

\begin{remark}[Folding for positive parameters]\label{rem:positive-folding}
For $m\in\mathbb N_0$ and
$\delta\in[m-\frac12,m+\frac12]$,
\begin{equation}\label{eq:triangular-q}
 q_\delta=(-1)^{m+1}(\delta-m).
\end{equation}
Thus, for $\delta\ge\frac12$, the exponent of the product determinant is
$-\operatorname{dist}(\delta,\mathbb Z)^2$.  Away from half-integers,
\[
 \det\bigl(I_N-(H_N^\delta)^2\bigr)
 \asymp_\delta N^{-\operatorname{dist}(\delta,\mathbb Z)^2}.
\]
The return to zero at nonnegative integers is unavoidable: for
$m\in\mathbb N_0$, the sine factor vanishes while all denominators are
nonzero, and hence $H_N^m=0$ exactly.  The two individual determinant
exponents are interchanged when $\delta$ is increased by one.
\end{remark}

The proof uses three established inputs, all stated explicitly below: Otte's
Hilbert determinant theorem, including its endpoint form; the spectral
description of the infinite generalized Hilbert matrix; and a classical
lower bound of order $(\log N)^{-2}$ for the gap between the norm of a finite
Hilbert matrix and $\pi$.  Away from $\mathcal C$, Otte's $O(1)$ remainder and
a trace-class comparison give an $O(1)$ remainder for the regularized
determinant.  Removing the regularization produces a finite-dimensional
factor.  The edge-state extraction lemma compares this factor, in Loewner
order, with the Gram matrix of the omitted eigenvector tails, and the latter
has an explicit Cauchy-matrix asymptotic.  At the negative critical
half-integers, an exact compression reduces the continuous part to the base
parameters $\alpha=\pm\frac12$; the closed spectral gap then costs only a
subpolynomial factor.

For $r=-\delta>\frac12$, let $m$ denote the number of edge eigenvectors.  The
algebraic core of the unfolding is
\begin{equation}\label{eq:phase-lift-intro}
 r=m+(-1)^m\frac1\pi\arcsin(\sin\pi r).
\end{equation}
The principal inverse sine gives the fractional part of the lifted phase,
whereas the number and parity of the edge eigenvectors give the missing branch
index.

\section{The infinite generalized Hilbert matrix}

Let $\ell^2(\Nzero)$ have standard orthonormal basis $(e_j)_{j\ge0}$.  For a
real parameter
$\alpha\notin\{0,-1,-2,\ldots\}$, let
\begin{equation}\label{eq:Halpha}
 \Hh_\alpha:=\left(\frac1{j+k+\alpha}\right)_{j,k\in\Nzero}
\end{equation}
be the bounded self-adjoint generalized Hilbert operator on
$\ell^2(\Nzero)$, and define its normalized form
\begin{equation}\label{eq:Aalpha}
 A_\alpha:=\frac{\sin(\pi\alpha)}\pi\Hh_\alpha.
\end{equation}
Let $P_N$ be the orthogonal projection onto
$\operatorname{span}\{e_0,\ldots,e_{N-1}\}$ and put $Q_N:=\id-P_N$.  With
$\alpha=1+\delta$,
\begin{equation}\label{eq:compression}
 H_N^\delta=P_NA_\alpha P_N\big|_{P_N\ell^2(\Nzero)}.
\end{equation}

Recall that
\[
 H^2(\D)
 :=\left\{f(z)=\sum_{n=0}^\infty a_nz^n:
              \sum_{n=0}^\infty|a_n|^2<\infty\right\},
 \qquad
 \|f\|_{H^2}^2:=\sum_{n=0}^\infty|a_n|^2.
\]
Thus the coefficient map
\[
 U:\ell^2(\Nzero)\longrightarrow H^2(\D),
 \qquad
 (Uu)(z)=\sum_{n=0}^\infty u(n)z^n,
\]
is unitary and identifies the entries of $u$ with the Taylor coefficients of
the analytic function $Uu$.  The same coefficients also have a Fourier
interpretation: the radial boundary values exist in $L^2(\mathbb T)$ and
satisfy
\[
 (Uu)^*(e^{it})=\sum_{n=0}^\infty u(n)e^{int}
 \quad\text{in }L^2(\mathbb T).
\]
Hence $u(n)$ are the nonnegative Fourier coefficients of the boundary
function; no pointwise convergence on $\mathbb T$ is being assumed.

The following proposition is the only spectral input.  It combines
Rosenblum's diagonalization with the explicit eigenfunction formulas in
\cite{RosenblumII,AlemanMontesSarafoleanu,Silbermann}.

\begin{proposition}[Spectral data]\label{prop:spectral-data}
Let $\alpha=1+\delta\in\R\setminus\mathbb Z$ and set
$s:=\sin(\pi\alpha)$.  Then:

\begin{enumerate}
\item The absolutely continuous spectrum of $A_\alpha$ is
\begin{equation}\label{eq:ac-spectrum}
 \operatorname{spec}_{\mathrm{ac}}(A_\alpha)
 =[\min\{0,s\},\max\{0,s\}],
\end{equation}
with multiplicity one.

\item The point spectrum of $A_\alpha$ is contained in $\{-1,+1\}$.  Define
\begin{align}
 J^-(\delta)
 &:=\left\{j\in\Nzero:j<-\delta-\frac12,\ j\ \text{even}\right\},
 \label{eq:Jminus}\\
 J^+(\delta)
 &:=\left\{j\in\Nzero:j<-\delta-\frac12,\ j\ \text{odd}\right\}.
 \label{eq:Jplus}
\end{align}
For each admissible $j$, let $u_j$ be the Taylor coefficient vector of
\begin{equation}\label{eq:fj}
 f_j(z):=(1-z)^{-\alpha-j}(1+z)^j.
\end{equation}
Then
\begin{equation}\label{eq:eigenvalue-fj}
 A_\alpha u_j=(-1)^j u_j,
\end{equation}
and
\begin{equation}\label{eq:kernel-basis}
 \Ker(\id\pm A_\alpha)
 =\operatorname{span}\{u_j:j\in J^\pm(\delta)\}.
\end{equation}

\item The operator $A_\alpha$ is a self-adjoint contraction.  If $|s|<1$,
then, after the kernels in \eqref{eq:kernel-basis} are removed, the spectra of
$\id\pm A_\alpha$ are bounded away from zero.
\end{enumerate}
\end{proposition}

The restriction $j<-\delta-\frac12$ in
\eqref{eq:Jminus}--\eqref{eq:Jplus} is exactly the condition that
$f_j\in H^2(\D)$.  Indeed, near $z=1$,
$f_j(z)$ is a nonzero constant times $(1-z)^{-\alpha-j}$, and this singularity
belongs to $H^2$ precisely when $\alpha+j<\frac12$.

At nonnegative integral $\delta$, the sine factor in
\eqref{eq:def-Hdelta} vanishes and $H_N^\delta=0$.  These elementary cases do
not require the spectral proposition.

\section{The principal-branch Hilbert contribution}

For $a>0$, write
\begin{equation}\label{eq:finite-Hilbert}
 \Hh_{N,a}:=\left(\frac1{j+k+a}\right)_{j,k=0}^{N-1}.
\end{equation}
We use the following form of Otte's generalized Hilbert determinant theorem.

\begin{proposition}[Hilbert determinant asymptotics]\label{prop:Hilbert-det}
Let $a\ge\frac12$ and $\beta\in[-1,1]$.  Set
\begin{equation}\label{eq:hilbert-exponent}
 h(\beta):=
 \frac{\arcsin^2\beta+\pi\arcsin\beta}{2\pi^2},
\end{equation}
where the inverse sine is the principal branch.  Then
\begin{equation}\label{eq:Otte-strong}
 \log\det\left(\id_N-\frac\beta\pi\Hh_{N,a}\right)
 =-h(\beta)\log N+
 \begin{cases}
  O_{a,\beta}(1),&\beta<1,\\
  o(\log N),&\beta=1.
 \end{cases}
\end{equation}
\end{proposition}

Fedele--Gebert first obtained the leading term for $|\beta|<1$
\cite{FedeleGebert}.  Otte proved the generalized-parameter statement and
improved the remainder to $O(1)$ for $\beta<1$; the endpoint $\beta=1$ is
proved with an $o(\log N)$ remainder \cite{Otte}.  For fixed $a$, Otte's more precise logarithmic scale
differs from a constant multiple of $\log N$ by $O(1)$, which gives
\eqref{eq:Otte-strong} in the form used here.

Suppose first that $\delta\ge-\frac12$.  Then
$\alpha=1+\delta\ge\frac12$ and
\[
 D_N^\pm(\delta)
 =\det\left(\id_N-\frac{\mp s_\delta}{\pi}\Hh_{N,\alpha}\right).
\]
Since $\arcsin(\mp s_\delta)=\mp\pi q_\delta$,
Proposition~\ref{prop:Hilbert-det} gives
\begin{equation}\label{eq:Otte-paired}
 \log D_N^\pm(\delta)
 =-\frac{q_\delta^2\mp q_\delta}{2}\log N+
 \begin{cases}
  O_\delta(1),&\mp s_\delta<1,\\
  o(\log N),&\mp s_\delta=1.
 \end{cases}
\end{equation}
If $\delta\notin\mathcal C$, then $|s_\delta|<1$ and both signs have
bounded remainders.  If $\delta\in\mathcal C$, one sign has endpoint
coupling $1$ and the other has coupling $-1$.  Since
$\theta_\delta=q_\delta$ throughout $\delta\ge-\frac12$, this proves
Theorem~\ref{thm:main} in that range.

For $\alpha<\frac12$, Otte's theorem does not apply directly.  The next two
lemmas show that, away from the threshold couplings $|s|=1$, the regularized
operator may nevertheless be compared with the standard Hilbert matrix.

\begin{lemma}[Trace-class shift]\label{lem:trace-class-shift}
If $\alpha,\beta\in\R\setminus\{0,-1,-2,\ldots\}$, then
\begin{equation}\label{eq:trace-class-difference}
 \Hh_\alpha-\Hh_\beta\in\cS.
\end{equation}
\end{lemma}

\begin{proof}
Choose $r\in\Nzero$ so large that $\alpha+2r>0$ and $\beta+2r>0$.
Deleting the first $r$ rows and columns changes the difference by a finite-rank
operator.  After reindexing, the remaining block is
$\Hh_{\alpha+2r}-\Hh_{\beta+2r}$.

Assume, without loss of generality, that
$a:=\alpha+2r<b:=\beta+2r$.  Its entries satisfy
\[
 \frac1{j+k+a}-\frac1{j+k+b}
 =\int_0^1 t^{j+k}\bigl(t^{a-1}-t^{b-1}\bigr)\,dt,
\]
so the difference is positive.  Its diagonal sum is finite:
\[
 \sum_{n=0}^\infty
 \left(\frac1{2n+a}-\frac1{2n+b}\right)<\infty.
\]
A positive bounded operator with finite diagonal sum is trace class.  The case
$b<a$ follows after changing the sign.
\end{proof}

\begin{lemma}[Stable trace-class perturbations]\label{lem:stable-traceclass}
Let $R,S$ be bounded self-adjoint operators on a Hilbert space, suppose that
$R\ge c\id$ and $S\ge c\id$ for some $c>0$, and assume $S-R\in\cS$.  Then
\begin{equation}\label{eq:stable-determinant}
 \left|
 \log\det(P_NSP_N)-\log\det(P_NRP_N)
 \right|
 \le c^{-1}\|S-R\|_1.
\end{equation}
The determinants are taken on the finite-dimensional space $P_N\ell^2$.
\end{lemma}

\begin{proof}
Set $R_N:=P_NRP_N$ and $K_N:=P_N(S-R)P_N$.  For $0\le t\le1$,
\[
 R_N+tK_N=(1-t)P_NRP_N+tP_NSP_N\ge cP_N.
\]
Therefore
\[
 \frac{d}{dt}\log\det(R_N+tK_N)
 =\Tr\bigl((R_N+tK_N)^{-1}K_N\bigr),
\]
and the absolute value of the right-hand side is bounded by
$c^{-1}\|K_N\|_1\le c^{-1}\|S-R\|_1$.  Integration over $t\in[0,1]$
proves the claim.
\end{proof}

Fix henceforth a noncritical parameter
\begin{equation}\label{eq:nonthreshold}
 |s|=|\sin(\pi\alpha)|<1.
\end{equation}
Define
\begin{equation}\label{eq:Tpm}
 T^\pm:=\id\pm A_\alpha,
 \qquad
 K^\pm:=\Ker T^\pm,
\end{equation}
let $\Pi^\pm$ be the orthogonal projection onto $K^\pm$, and set
\begin{equation}\label{eq:Ttilde-pm}
 \widetilde T^{\,\pm}:=T^\pm+\Pi^\pm.
\end{equation}
By Proposition~\ref{prop:spectral-data}, the two regularized operators are
bounded below by positive constants.  The same is true of
\begin{equation}\label{eq:Rpm}
 R^\pm:=\id\pm\frac{s}{\pi}\Hh_1,
\end{equation}
because $\pi^{-1}\Hh_1$ has spectrum $[0,1]$ and $|s|<1$.  Moreover,
\begin{equation}\label{eq:regularized-difference}
 \widetilde T^{\,\pm}-R^\pm
 =\pm\frac{s}{\pi}(\Hh_\alpha-\Hh_1)+\Pi^\pm
 \in\cS.
\end{equation}
Lemmas~\ref{lem:trace-class-shift} and \ref{lem:stable-traceclass}, together
with Proposition~\ref{prop:Hilbert-det} applied at $a=1$, give, with
\begin{equation}\label{eq:def-q}
 q:=\frac1\pi\arcsin(s),
\end{equation}
the regularized asymptotics
\begin{equation}\label{eq:regularized-asymptotic}
 \log\det(P_N\widetilde T^{\,\pm}P_N)
 =-\frac{q^2\mp q}{2}\log N+O_\delta(1).
\end{equation}
This is the contribution that depends only on the principal inverse sine.

\section{Extracting the edge eigenspaces}

The next lemma converts the exact kernel of an infinite operator into the small
finite-dimensional factor in its finite-section determinant.  The comparison
in \eqref{eq:loewner-tail} is the key point: it says that this factor is,
up to fixed constants, the Gram matrix of the omitted tails of the kernel
vectors.

\begin{lemma}[Edge-eigenspace extraction]\label{lem:edge-extraction}
Let $T\ge0$ be bounded on $\ell^2(\Nzero)$.  Suppose that
$K:=\Ker T$ has finite dimension $m$ and that $T$ is bounded below by a
positive constant on $K^\perp$.  Let $\Pi$ be the orthogonal projection onto
$K$, set $\widetilde T:=T+\Pi$, and let $V:\C^m\to K$ be an isometry.  Then
\begin{equation}\label{eq:edge-factorization}
 \det(P_NTP_N)
 =\det(P_N\widetilde TP_N)\det M_N,
\end{equation}
where
\begin{equation}\label{eq:def-MN}
 M_N:=\id_m-V^*P_N(P_N\widetilde TP_N)^{-1}P_NV.
\end{equation}
There are constants $0<c\le C<\infty$, independent of $N$, such that
\begin{equation}\label{eq:loewner-tail}
 c\,V^*Q_NV\le M_N\le C\,V^*Q_NV
\end{equation}
in the Loewner order on $m\times m$ matrices.  More precisely, one may take
\begin{equation}\label{eq:loewner-constants}
 c=\|\widetilde T^{-1/2}\|^{-2},
 \qquad
 C=\|\widetilde T^{1/2}\|^2.
\end{equation}
If $V^*Q_NV$ is positive definite, then
\begin{equation}\label{eq:edge-logdet}
 \log\det M_N
 =\log\det(V^*Q_NV)+O(1).
\end{equation}
\end{lemma}

\begin{proof}
Since $\Pi=VV^*$,
\[
 P_NTP_N=P_N\widetilde TP_N-(P_NV)(P_NV)^*.
\]
The matrix determinant lemma gives
\eqref{eq:edge-factorization}--\eqref{eq:def-MN}.

Put $B:=\widetilde T^{1/2}$.  The hypotheses imply that $B$ is boundedly
invertible.  Moreover, $T=0$ and $\Pi=\id$ on $K$, so
\begin{equation}\label{eq:B-on-kernel}
 BV=V.
\end{equation}
Let
\[
 C_N:=BP_N:P_N\ell^2\longrightarrow\ell^2.
\]
Then
\[
 C_N^*C_N=P_N\widetilde TP_N,
 \qquad
 C_N^*V=P_NBV=P_NV.
\]
Because $C_N$ is injective and has finite-dimensional range,
\begin{equation}\label{eq:EN-projection}
 E_N:=C_N(C_N^*C_N)^{-1}C_N^*
\end{equation}
is the orthogonal projection onto
$\Ran C_N=BP_N\ell^2$.  Hence, for every $x\in\C^m$,
\begin{align}
 x^*M_Nx
 &=\|Vx\|^2-
   \bigl\langle C_N^*Vx,(C_N^*C_N)^{-1}C_N^*Vx\bigr\rangle\notag\\
 &=\langle Vx,(\id-E_N)Vx\rangle\notag\\
 &=\|(\id-E_N)Vx\|^2\notag\\
 &=\dist\bigl(Vx,BP_N\ell^2\bigr)^2.
 \label{eq:MN-distance}
\end{align}
This is the geometric meaning of $M_N$.

Using \eqref{eq:B-on-kernel},
\begin{align*}
 \dist\bigl(Vx,BP_N\ell^2\bigr)
 &=\inf_{y\in P_N\ell^2}\|Vx-By\|\\
 &=\inf_{y\in P_N\ell^2}\|B(Vx-y)\|.
\end{align*}
For every $z\in\ell^2$,
\[
 \|B^{-1}\|^{-1}\|z\|
 \le\|Bz\|
 \le\|B\|\|z\|.
\]
Taking the infimum over $y\in P_N\ell^2$ and using that $P_N$ is an
orthogonal projection gives
\begin{equation}\label{eq:distance-comparison}
 \|B^{-1}\|^{-1}\|Q_NVx\|
 \le\dist\bigl(Vx,BP_N\ell^2\bigr)
 \le\|B\|\|Q_NVx\|.
\end{equation}
Since
\[
 \|Q_NVx\|^2=x^*V^*Q_NVx,
\]
combining \eqref{eq:MN-distance} and \eqref{eq:distance-comparison} proves
\eqref{eq:loewner-tail} with the constants in
\eqref{eq:loewner-constants}.

Finally, put $G_N:=V^*Q_NV>0$.  Conjugating
\eqref{eq:loewner-tail} by $G_N^{-1/2}$ gives
\[
 c\id_m\le G_N^{-1/2}M_NG_N^{-1/2}\le C\id_m.
\]
Therefore
\[
 c^m\le\frac{\det M_N}{\det G_N}\le C^m.
\]
Since $m$ is fixed, this is equivalent to \eqref{eq:edge-logdet}.
\end{proof}

\begin{remark}[Meaning of the Loewner comparison]
The matrix $M_N$ measures the squared distance of the kernel vectors $Vx$ from
the distorted finite-section space $\widetilde T^{1/2}P_N\ell^2$, whereas
$V^*Q_NV$ measures their ordinary squared tails.  The boundedly invertible map
$\widetilde T^{1/2}$ changes these distances by at most its condition-number
constants, which is exactly \eqref{eq:loewner-tail}.  In the special case
$\widetilde T=\id$, there is no distortion and
\[
 M_N=\id_m-V^*P_NV=V^*Q_NV
\]
holds identically.
\end{remark}

Applying Lemma~\ref{lem:edge-extraction} to $T^\pm$ reduces the remaining
problem to the tails of the explicit vectors $u_j$ from
Proposition~\ref{prop:spectral-data}.

\section{The tail Gram determinant}

Let $u_j(n)$ denote the $n$th Taylor coefficient of $f_j$ in
\eqref{eq:fj}.

\begin{lemma}[Eigenvector tails]\label{lem:eigenvector-tails}
Assume $\alpha\notin\mathbb Z$ and let $j\in\Nzero$ satisfy
$\alpha+j<\frac12$.  Then
\begin{equation}\label{eq:coeff-asymptotic}
 u_j(n)
 =c_j n^{\alpha+j-1}\bigl(1+O(n^{-1})\bigr),
 \qquad
 c_j:=\frac{2^j}{\Gamma(\alpha+j)}\ne0.
\end{equation}
If $J\subset\{j\in\Nzero:\alpha+j<\frac12\}$ is finite, then
\begin{equation}\label{eq:tail-gram-asymptotic}
 \det\left(\sum_{n=N}^\infty u_i(n)u_j(n)\right)_{i,j\in J}
 =C_J
 N^{\sum_{j\in J}(2\delta+2j+1)}\bigl(1+o(1)\bigr)
\end{equation}
for some constant $C_J>0$.
\end{lemma}

\begin{proof}
Since
\[
 (1+z)^j=\sum_{\ell=0}^j\binom j\ell z^\ell,
\]
the binomial series for $(1-z)^{-\alpha-j}$ gives, for $n\ge j$,
\[
 u_j(n)
 =\sum_{\ell=0}^j\binom j\ell
 \frac{\Gamma(n-\ell+\alpha+j)}
 {\Gamma(\alpha+j)\Gamma(n-\ell+1)}.
\]
The gamma-ratio asymptotic
\[
 \frac{\Gamma(n+a)}{\Gamma(n+b)}
 =n^{a-b}\bigl(1+O(n^{-1})\bigr)
\]
then yields \eqref{eq:coeff-asymptotic}, because
$\sum_{\ell=0}^j\binom j\ell=2^j$.

For $i,j\in J$, the hypotheses imply $2\delta+i+j<-1$.  Therefore
\begin{equation}\label{eq:tail-pairing}
 \sum_{n=N}^\infty u_i(n)u_j(n)
 =\frac{c_ic_j}{-(2\delta+i+j+1)}
 N^{2\delta+i+j+1}\bigl(1+o(1)\bigr).
\end{equation}
Factor $N^{\delta+i+1/2}$ from row $i$ and
$N^{\delta+j+1/2}$ from column $j$.  The remaining matrix converges to
\begin{equation}\label{eq:Cauchy-limit}
 \left(\frac{c_ic_j}{-(2\delta+i+j+1)}\right)_{i,j\in J}.
\end{equation}
This limit is positive definite: it is the Gram matrix of the linearly
independent functions $c_jx^{\delta+j}$ in $L^2(1,\infty)$.  Its determinant
is therefore positive, proving \eqref{eq:tail-gram-asymptotic}.
\end{proof}

For either sign, changing from the basis
$\{u_j:j\in J^\pm(\delta)\}$ to an orthonormal basis of $K^\pm$ multiplies the
tail Gram determinant by a fixed positive constant.  Lemmas
\ref{lem:edge-extraction} and \ref{lem:eigenvector-tails} therefore give
\begin{equation}\label{eq:edge-contribution}
 \log\det M_N^\pm
 =\left(
 \sum_{j\in J^\pm(\delta)}(2\delta+2j+1)
 \right)\log N+O(1),
\end{equation}
where $M_N^\pm$ is the matrix in \eqref{eq:def-MN} for $T=T^\pm$.
Combining \eqref{eq:regularized-asymptotic} and
\eqref{eq:edge-contribution} yields
\begin{equation}\label{eq:master-formula}
 \log D_N^\pm(\delta)
 =\left[
 -\frac{q^2\mp q}{2}
 +\sum_{j\in J^\pm(\delta)}(2\delta+2j+1)
 \right]\log N+O_\delta(1)
\end{equation}
whenever $|\sin(\pi(1+\delta))|<1$.

Formula~\eqref{eq:master-formula} is the unsimplified strong asymptotic.  Its
first term depends only on the principal inverse sine.  Its finite sum is the
contribution of the exact eigenvalues at $\mp1$ to
$D_N^\pm=\det(\id_N\pm H_N^\delta)$.

\section{Recovery of the original sine argument}

Assume that $\delta<-\frac12$ and $\delta\notin\mathcal C$.  Put
\begin{equation}\label{eq:r-m-q}
 r:=-\delta>\frac12,
 \qquad
 m:=\#\{j\in\Nzero:j<r-\tfrac12\},
 \qquad
 q:=\frac1\pi\arcsin(\sin\pi r).
\end{equation}
More generally, for every $r>\frac12$ the above definition of $m$ gives
$r\in(m-\frac12,m+\frac12]$, and on this interval
\begin{equation}\label{eq:lift}
 q=(-1)^m(r-m),
 \qquad\text{equivalently}\qquad
 r=m+(-1)^mq.
\end{equation}
The $m$ edge eigenvectors are indexed by $j=0,\ldots,m-1$, with
alternating eigenvalues.  Define
\begin{equation}\label{eq:Jm-pm}
 J_m^+:=\{0\le j<m:j\ \text{odd}\},
 \qquad
 J_m^-:=\{0\le j<m:j\ \text{even}\}.
\end{equation}
Thus $J^\pm(\delta)=J_m^\pm$.

\begin{lemma}[Unfolding identity]\label{lem:unfolding}
With $r,m,q$ as in \eqref{eq:r-m-q}--\eqref{eq:lift},
\begin{equation}\label{eq:unfolding-identity}
 -\frac{q^2\mp q}{2}
 +\sum_{j\in J_m^\pm}(1-2r+2j)
 =-\frac{r^2\mp r}{2}.
\end{equation}
\end{lemma}

\begin{proof}
Suppose first that $m=2a$.  Then $q=r-2a$, and
\begin{align*}
 \sum_{j\in J_m^-}(1-2r+2j)
 &=a(1-2r)+2a(a-1),\\
 \sum_{j\in J_m^+}(1-2r+2j)
 &=a(1-2r)+2a^2.
\end{align*}
Substitution gives \eqref{eq:unfolding-identity} for both signs.

If $m=2a+1$, then $q=2a+1-r$, and
\begin{align*}
 \sum_{j\in J_m^-}(1-2r+2j)
 &=(a+1)(1-2r)+2a(a+1),\\
 \sum_{j\in J_m^+}(1-2r+2j)
 &=a(1-2r)+2a^2.
\end{align*}
The same substitution proves the identity.
\end{proof}

Since $2\delta+2j+1=1-2r+2j$, Lemma~\ref{lem:unfolding} turns
\eqref{eq:master-formula} into
\begin{equation}\label{eq:unfolded-final}
 \log D_N^\pm(\delta)
 =-\frac{r^2\mp r}{2}\log N+O_\delta(1)
 =-\frac{\delta^2\pm\delta}{2}\log N+O_\delta(1).
\end{equation}
This proves the strong statement in Theorem~\ref{thm:main} for
$\delta<-\frac12$ away from $\mathcal C$.  Since
\begin{equation}\label{eq:product-determinant}
 \det\bigl(\id_N-(H_N^\delta)^2\bigr)
 =D_N^+(\delta)D_N^-(\delta),
\end{equation}
the two powers add to $-\delta^2$.  The negative critical cases are completed
in the next section.

The recovery of the sine argument is now explicit.  The absolutely continuous
part of the operator depends on $\delta$ only through
$s_\delta=\sin(-\pi\delta)$ and therefore yields the principal value
$q=\pi^{-1}\arcsin(s_\delta)$.  Whenever $r=-\delta$ crosses a half-integer,
a new function $f_j$ enters $H^2(\D)$ and produces an eigenvalue at $+1$ or
$-1$.  The tail of that eigenvector contributes the additional power of $N$
required to pass to the next branch.  Identity~\eqref{eq:lift} states this
precisely: the continuous spectral contribution and the edge-eigenvector
count reconstruct the lifted phase $r$, or equivalently the original sine
argument $-\pi\delta$.

For $-\frac12\le\delta\le\frac12$, the principal branch already satisfies
$q_\delta=-\delta$, so the same linear formula continues without a discrete
correction.  For $\delta>\frac12$, however, Proposition~\ref{prop:spectral-data}
has no edge eigenvectors: the generalized Hilbert parameter
$\alpha=1+\delta$ lies above the point-spectrum range.  The determinant
therefore retains the principal phase $q_\delta$.  In particular,
\[
 |q_\delta|=\operatorname{dist}(\delta,\mathbb Z),
 \qquad
 q_{\delta+1}=-q_\delta.
\]
The squared-determinant exponent is consequently $1$-periodic on this side,
and the two individual exponents are interchanged by $\delta\mapsto\delta+1$.
The exact identity $H_N^n=0$ at every nonnegative integer $n$ makes this
folding unavoidable.

\section{Negative critical half-integers}\label{sec:threshold}

It remains to treat
\begin{equation}\label{eq:threshold-parameter}
 \delta=-m-\frac12,
 \qquad m=1,2,\ldots.
\end{equation}
Then
\[
 \alpha=1+\delta=\frac12-m,
 \qquad
 \sin(\pi\alpha)=(-1)^m.
\]
It is convenient in this section to write
\begin{equation}\label{eq:KTa}
 K_a:=\frac1\pi\Hh_a,
 \qquad
 T_a:=\id-K_a,
 \qquad
 T_{N,a}:=P_NT_aP_N\big|_{P_N\ell^2(\Nzero)}.
\end{equation}
Thus $A_\alpha=(-1)^mK_\alpha$.  The determinant whose absolutely continuous
spectrum reaches zero is $D_N^-$ if $m$ is even and $D_N^+$ if $m$ is odd; in
either case it equals $\det T_{N,\alpha}$.

The proof requires a weak quantitative form of endpoint coercivity.  For a
sequence $0<\gamma_N\le1$, we write $\gamma_N=N^{-o(1)}$ when
\begin{equation}\label{eq:subpolynomial-definition}
 -\log\gamma_N=o(\log N),
\end{equation}
equivalently, $\gamma_N^{-1}=N^{o(1)}$.  Such a loss is negligible on the
scale of powers of $N$ but need not be bounded.

\begin{lemma}[Endpoint coercivity at $a=\frac12$]
\label{lem:positive-endpoint-coercivity}
There is a sequence $\gamma_N=N^{-o(1)}$ such that
\begin{equation}\label{eq:positive-endpoint-coercivity}
 T_{N,1/2}\ge\gamma_N\id_N.
\end{equation}
\end{lemma}

\begin{proof}
Write $T_{N,1/2}$ in the decomposition
$\C e_0\oplus\operatorname{span}\{e_1,\ldots,e_{N-1}\}$ as
\begin{equation}\label{eq:block-positive-endpoint}
 T_{N,1/2}
 =\begin{pmatrix} a_N&b_N^*\\ b_N&R_N\end{pmatrix},
 \qquad
 R_N=T_{N-1,5/2}.
\end{equation}
The moment representation
\[
 \Hh_{n,a}=\int_0^1 x^{a-1}
 (1,x,\ldots,x^{n-1})^{\mathsf T}
 (1,x,\ldots,x^{n-1})\,dx
\]
shows that $\Hh_{n,5/2}\le\Hh_{n,2}$.  The classical finite Hilbert-matrix
estimate
\begin{equation}\label{eq:Wilf-gap}
 \pi-\|\Hh_{n,2}\|\ge \frac{c}{\log^2(n+1)}
\end{equation}
for all sufficiently large $n$ therefore gives
\begin{equation}\label{eq:R-gap}
 R_N\ge\frac{c}{\log^2(N+1)}\id_{N-1}.
\end{equation}
The sharper asymptotic behind \eqref{eq:Wilf-gap} is due to de Bruijn and
Wilf; only this lower bound is needed here \cite{Wilf}.

Let
\[
 \rho_N:=a_N-b_N^*R_N^{-1}b_N
 =\frac{\det T_{N,1/2}}{\det T_{N-1,5/2}}
\]
be the scalar Schur complement.  Both determinants are positive.  Otte's
endpoint theorem, applied once with parameter $1/2$ and once with parameter
$5/2$, gives
\begin{equation}\label{eq:rho-subpoly}
 \log\rho_N=o(\log N),
 \qquad
 \rho_N^{-1}=N^{o(1)}.
\end{equation}
The second assertion follows from the first because $\rho_N>0$.

Moreover, $\|b_N\|$ is bounded uniformly in $N$.  Put
\[
 U_N:=\begin{pmatrix}1&0\\R_N^{-1}b_N&\id\end{pmatrix}.
\]
The Schur factorization may be written as the congruence
\begin{equation}\label{eq:endpoint-LDU}
 T_{N,1/2}
 =U_N^*
   \begin{pmatrix}\rho_N&0\\0&R_N\end{pmatrix}
  U_N.
\end{equation}
By \eqref{eq:R-gap},
$\|R_N^{-1}b_N\|=O(\log^2(N+1))$, and hence
\begin{equation}\label{eq:endpoint-U-bound}
 \|U_N\|+\|U_N^{-1}\|=O(\log^2(N+1)).
\end{equation}
For every vector $z$,
\[
 \langle z,T_{N,1/2}z\rangle
 \ge
 \min\left\{\rho_N,\frac{c}{\log^2(N+1)}\right\}
 \|U_Nz\|^2.
\]
Using $\|U_Nz\|\ge\|U_N^{-1}\|^{-1}\|z\|$ and
\eqref{eq:rho-subpoly}--\eqref{eq:endpoint-U-bound}, the right-hand side is
bounded below by $\gamma_N\|z\|^2$ for a sequence
$\gamma_N=N^{-o(1)}$.  This proves
\eqref{eq:positive-endpoint-coercivity}.
\end{proof}

The other base parameter, $a=-\frac12$, is not covered by Otte's theorem.  It
can nevertheless be reduced to the two positive parameters $a=\frac12$ and
$a=\frac32$ by elementary eigenvalue interlacing.

\begin{proposition}[The base negative threshold]\label{prop:negative-half-base}
As $N\to\infty$,
\begin{equation}\label{eq:negative-half-det}
 \log\det T_{N,-1/2}=-\frac38\log N+o(\log N).
\end{equation}
Moreover, there is a sequence $\widetilde\gamma_N=N^{-o(1)}$ such that
\begin{equation}\label{eq:negative-half-coercivity}
 T_{N,-1/2}\ge\widetilde\gamma_N\id_N.
\end{equation}
\end{proposition}

\begin{proof}
Set $L_N:=K_{N,-1/2}$.  Its lower-right principal submatrix is
$K_{N-1,3/2}>0$, whereas
$\langle e_0,L_Ne_0\rangle=-2/\pi<0$.  Cauchy interlacing therefore shows that
$L_N$ has exactly one negative eigenvalue.  Write its eigenvalues as
\begin{equation}\label{eq:L-eigenvalues}
 \lambda_1\ge\cdots\ge\lambda_{N-1}>0>\lambda_N.
\end{equation}

The infinite operator $K_{-1/2}$ has the eigenvalue $-1$ with eigenvector given
by the Taylor coefficients of
\begin{equation}\label{eq:negative-half-eigenvector}
 (1-z)^{1/2}.
\end{equation}
These coefficients are $O(n^{-3/2})$.  If $u$ is the normalized coefficient
vector, then $\|Q_Nu\|=O(N^{-1})$, and hence
\[
 \|(L_N+\id_N)P_Nu\|
 =\|P_NK_{-1/2}Q_Nu\|=O(N^{-1}).
\]
Since $\|P_Nu\|\to1$, the spectral theorem implies that the distance from
$-1$ to the spectrum of $L_N$ is $O(N^{-1})$.  The only negative eigenvalue is
$\lambda_N$, while all other eigenvalues are positive.  Hence
\begin{equation}\label{eq:negative-eigenvalue-rate}
 \lambda_N=-1+O(N^{-1}).
\end{equation}

Delete the first row of $L_N$, and denote the resulting
$(N-1)\times N$ matrix by $B_N$.  After reindexing, $B_N$ is the first
$N-1$ rows of $K_{N,1/2}$.  For $j=1,\ldots,N-1$, singular-value
interlacing under deletion of one row and the ideal property give
\begin{equation}\label{eq:singular-interlace-full}
 \sigma_{j+1}(L_N)
 \le \sigma_j(B_N)
 \le \sigma_j(K_{N,1/2}).
\end{equation}
In particular,
\begin{equation}\label{eq:singular-interlace-base}
 \sigma_2(L_N)
 \le\|K_{N,1/2}\|
 \le1-\gamma_N,
\end{equation}
where $\gamma_N=N^{-o(1)}$ is supplied by
Lemma~\ref{lem:positive-endpoint-coercivity}.  Since
$N^{-1}=o(\gamma_N)$, equations \eqref{eq:negative-eigenvalue-rate} and
\eqref{eq:singular-interlace-base} imply that $|\lambda_N|$ is the largest
singular value of $L_N$ for all sufficiently large $N$.

Let $\mu_1\ge\cdots\ge\mu_N>0$ be the eigenvalues of $K_{N,1/2}$.  Once
$|\lambda_N|$ is the largest singular value, the remaining singular values of
$L_N$ are $\lambda_1,\ldots,\lambda_{N-1}$ in that order.  Thus
\eqref{eq:singular-interlace-full} gives the explicit comparison
\begin{equation}\label{eq:positive-eigenvalue-comparison}
 \lambda_j=\sigma_{j+1}(L_N)
 \le\sigma_j(B_N)
 \le\mu_j,
 \qquad j=1,\ldots,N-1.
\end{equation}
The normalized infinite operator $K_{-1/2}$ is a contraction, and so is its
compression $L_N$; in particular, $-1\le\lambda_N<0$.  Consequently,
\begin{align}
 \det(\id_N-L_N)
 &=(1-\lambda_N)\prod_{j=1}^{N-1}(1-\lambda_j)\notag\\
 &\ge\prod_{j=1}^{N-1}(1-\mu_j)
 \ge\det(\id_N-K_{N,1/2}).
 \label{eq:base-lower-bound}
\end{align}
For the reverse estimate, let
$\nu_1\ge\cdots\ge\nu_{N-1}>0$ be the eigenvalues of the lower-right
principal submatrix $K_{N-1,3/2}$.  Cauchy interlacing gives
$\lambda_j\ge\nu_j$ for $j=1,\ldots,N-1$.  Since
$1-\lambda_N\le2$, it follows that
\begin{equation}\label{eq:base-upper-bound}
 \det(\id_N-L_N)
 \le2\prod_{j=1}^{N-1}(1-\nu_j)
 =2\det(\id_{N-1}-K_{N-1,3/2}).
\end{equation}
Otte's endpoint theorem applied to $a=\frac12$ and $a=\frac32$ in
\eqref{eq:base-lower-bound}--\eqref{eq:base-upper-bound} proves
\eqref{eq:negative-half-det}.

Finally, \eqref{eq:positive-eigenvalue-comparison} and
Lemma~\ref{lem:positive-endpoint-coercivity} imply
$\lambda_1\le1-\gamma_N$.  Since $1-\lambda_N>1$, this proves
\eqref{eq:negative-half-coercivity}.
\end{proof}

We next replace the uniform coercivity in
Lemma~\ref{lem:edge-extraction} by the subpolynomial estimates just obtained.
The following formulation is adapted to the Schur complements used below.

\begin{lemma}[Threshold tail extraction]\label{lem:threshold-tail-extraction}
Let $B\ge0$ be a bounded injective operator on $\ell^2(\Nzero)$ and suppose
that
\begin{equation}\label{eq:threshold-B-coercivity}
 P_MB P_M\ge\gamma_M P_M,
 \qquad
 \gamma_M=M^{-o(1)}.
\end{equation}
Let $L:\C^d\to\ell^2(\Nzero)$ be injective, put
$G_M:=L^*Q_ML$, and assume that $G_M$ is positive definite for all sufficiently
large $M$.  Suppose further that there are fixed constants $A>1$, $c>0$, and
$C<\infty$ such that
\begin{equation}\label{eq:scale-separation}
 L^*Q_{\lfloor M^A\rfloor}L
 \le C M^{-c}G_M
\end{equation}
for all sufficiently large $M$.  Define a positive $d\times d$ matrix by
\begin{equation}\label{eq:threshold-S-def}
 x^*S_Mx
 :=\inf_{y\in P_M\ell^2}
 \langle y-Lx,B(y-Lx)\rangle.
\end{equation}
Then
\begin{equation}\label{eq:threshold-S-asymptotic}
 \log\det S_M=\log\det G_M+o(\log M).
\end{equation}
\end{lemma}

\begin{proof}
Choosing $y=P_MLx$ in \eqref{eq:threshold-S-def} gives
\begin{equation}\label{eq:threshold-upper}
 S_M\le\|B\|G_M.
\end{equation}
For the lower bound, fix $y\in P_M\ell^2$, put
$w:=y-Lx$, and let $M_A:=\lfloor M^A\rfloor$.  By the triangle inequality,
\begin{align*}
 \|B^{1/2}w\|
 &\ge \|B^{1/2}P_{M_A}w\|
      -\|B^{1/2}Q_{M_A}w\|\\
 &\ge \gamma_{M_A}^{1/2}\|P_{M_A}w\|
      -\|B\|^{1/2}\|Q_{M_A}Lx\|.
\end{align*}
Because $y$ is supported in the first $M$ coordinates,
\[
 \|P_{M_A}w\|
 \ge\|(P_{M_A}-P_M)Lx\|.
\]
Condition \eqref{eq:scale-separation} therefore implies, uniformly in $x$,
\[
 \|(P_{M_A}-P_M)Lx\|^2
 \ge(1-CM^{-c})x^*G_Mx,
 \qquad
 \|Q_{M_A}Lx\|^2\le CM^{-c}x^*G_Mx.
\]
Since $\gamma_{M_A}=M^{-o(1)}$, one has
$M^{-c/2}\gamma_{M_A}^{-1/2}=M^{-c/2+o(1)}\to0$.  Hence the second term in
the lower bound for $\|B^{1/2}w\|$ is eventually at most one half of the
first, uniformly in $x$.  Taking the infimum over $y$ yields
\begin{equation}\label{eq:threshold-lower}
 S_M\ge\frac14\gamma_{M_A}G_M
\end{equation}
for all sufficiently large $M$.  Since $d$ is fixed,
\eqref{eq:threshold-upper} and \eqref{eq:threshold-lower} imply
\[
 d\log\gamma_{M_A}-d\log4
 \le \log\det S_M-\log\det G_M
 \le d\log\|B\|.
\]
Finally, $\log\gamma_{M_A}=o(\log M_A)=o(\log M)$, proving
\eqref{eq:threshold-S-asymptotic}.
\end{proof}

For the vectors in Lemma~\ref{lem:eigenvector-tails}, condition
\eqref{eq:scale-separation} follows directly from the Cauchy-Gram asymptotic.
Indeed, if the relevant indices are $j_1<\cdots<j_d$ and
\[
 D_M:=\operatorname{diag}
 \bigl(M^{\delta+j_1+1/2},\ldots,M^{\delta+j_d+1/2}\bigr),
\]
then \eqref{eq:tail-pairing} gives
\[
 G_M=D_M(C+o(1))D_M
\]
with $C>0$.  At $\lfloor M^A\rfloor$ the corresponding diagonal matrix is
$D_ME_M(\id+o(1))$, where
\[
 E_M:=\operatorname{diag}
 \bigl(M^{(A-1)(\delta+j_1+1/2)},\ldots,
       M^{(A-1)(\delta+j_d+1/2)}\bigr).
\]
Every exponent in $E_M$ is strictly negative.  Thus there is $c>0$ such
that $\|E_M\|^2\le M^{-c}$.  Since $C$ is positive definite, for all large
$M$ there are constants $0<c_0\le C_0<\infty$ with
\[
 c_0D_M^2\le G_M\le C_0D_M^2
\]
and
\[
 G_{\lfloor M^A\rfloor}
 \le C_0D_ME_M^2D_M
 \le \frac{C_0}{c_0}M^{-c}G_M.
\]
This is \eqref{eq:scale-separation}.  A fixed change of basis conjugates both
Gram matrices by the same invertible matrix and therefore does not affect the
conclusion.

\begin{proposition}[Threshold determinant asymptotics]
\label{prop:threshold-asymptotics}
Let $\delta=-m-\frac12$ with $m=1,2,\ldots$, and set
\begin{equation}\label{eq:threshold-q}
 q:=\frac1\pi\arcsin(\sin(-\pi\delta))=\frac{(-1)^m}{2}.
\end{equation}
Then
\begin{equation}\label{eq:threshold-master}
 \log D_N^\pm(\delta)
 =\left[
 -\frac{q^2\mp q}{2}
 +\sum_{j\in J^\pm(\delta)}(2\delta+2j+1)
 \right]\log N+o(\log N).
\end{equation}
Consequently,
\begin{equation}\label{eq:threshold-unfolded}
 \log D_N^\pm(\delta)
 =-\frac{\delta^2\pm\delta}{2}\log N+o(\log N),
\end{equation}
or equivalently
\begin{equation}\label{eq:threshold-power-form}
 D_N^\pm(\delta)
 =N^{-(\delta^2\pm\delta)/2+o(1)}.
\end{equation}
\end{proposition}

\begin{proof}
Put $\alpha=\frac12-m$ and $d:=\lfloor m/2\rfloor$.  Compression to the
coordinates $d,d+1,\ldots$ shifts the generalized Hilbert parameter by $2d$:
\begin{equation}\label{eq:threshold-compression}
 Q_dK_\alpha Q_d\cong K_a,
 \qquad
 a:=\alpha+2d=
 \begin{cases}
  \frac12,&m\ \text{even},\\
  -\frac12,&m\ \text{odd}.
 \end{cases}
\end{equation}
The kernel of $T_\alpha=\id-K_\alpha$ has dimension $d$ and is spanned by the
vectors $u_j$ with
\begin{equation}\label{eq:critical-J}
 0\le j<m,
 \qquad
 j\equiv m\pmod2.
\end{equation}
Indeed, $K_\alpha=(-1)^mA_\alpha$ and
$A_\alpha u_j=(-1)^ju_j$.

Let $E_d:=P_d\ell^2$ and $F_d:=Q_d\ell^2$.  The restriction from
$\Ker T_\alpha$ to $E_d$ is injective: a kernel vector vanishing on $E_d$
would, by \eqref{eq:threshold-compression}, give a vector in $\Ker T_a$, but
$T_{1/2}$ and $T_{-1/2}$ are injective.  Since both spaces have dimension $d$,
the restriction is an isomorphism.  Hence there is a map
$L:E_d\to F_d$ such that
\begin{equation}\label{eq:kernel-graph}
 \Ker T_\alpha=\{x\oplus Lx:x\in E_d\}.
\end{equation}

With respect to $E_d\oplus F_d$, write
\begin{equation}\label{eq:threshold-block-operator}
 T_\alpha=
 \begin{pmatrix}A&C^*\\ C&T_a\end{pmatrix}.
\end{equation}
Since every $x\oplus Lx$ lies in the kernel, the two block equations are
\[
 Cx+T_aLx=0,
 \qquad
 Ax+C^*Lx=0.
\]
Thus $C=-T_aL$ and, using self-adjointness, $A=L^*T_aL$.  Consequently,
\begin{equation}\label{eq:threshold-square-completion}
 \left\langle x\oplus y,T_\alpha(x\oplus y)\right\rangle
 =\langle y-Lx,T_a(y-Lx)\rangle.
\end{equation}

Let $M=N-d$.  The lower-right block of the $N$th finite section is
$T_{M,a}$ and is positive definite.  Taking its Schur complement gives
\begin{equation}\label{eq:threshold-det-factorization}
 \det T_{N,\alpha}=\det T_{M,a}\det S_M,
\end{equation}
where the Schur complement has the equivalent variational description
\begin{equation}\label{eq:threshold-variational}
 x^*S_Mx
 =\inf_{y\in P_MF_d}
 \langle y-Lx,T_a(y-Lx)\rangle.
\end{equation}
Proposition~\ref{prop:negative-half-base} and
Lemma~\ref{lem:positive-endpoint-coercivity} allow us to apply
Lemma~\ref{lem:threshold-tail-extraction} to \eqref{eq:threshold-variational}.
The columns of $L$ are a fixed basis of the kernel in
\eqref{eq:critical-J}; Lemma~\ref{lem:eigenvector-tails} is unchanged by this
fixed change of basis and by deleting finitely many coordinates.  Therefore
\begin{equation}\label{eq:critical-schur-exponent}
 \log\det S_M
 =\left(
   \sum_{\substack{0\le j<m\\j\equiv m\ (2)}}
   (2\delta+2j+1)
  \right)\log M+o(\log M).
\end{equation}
Both base determinants in \eqref{eq:threshold-compression} satisfy
\begin{equation}\label{eq:base-three-eighths}
 \log\det T_{M,a}=-\frac38\log M+o(\log M)
\end{equation}
by Otte's theorem when $a=\frac12$ and by
Proposition~\ref{prop:negative-half-base} when $a=-\frac12$.  This proves
\eqref{eq:threshold-master} for the critical determinant, because its
principal-branch term is $-3/8$.

For the other determinant the absolutely continuous part of
$\id+K_\alpha$ is bounded below by the identity.  After adding the projection
onto its finite-dimensional kernel at $-1$, the trace-class comparison and
Lemma~\ref{lem:edge-extraction} apply exactly as in Sections~3--5, now with the
reference operator $\id+\pi^{-1}\Hh_1$.  Proposition~\ref{prop:Hilbert-det}
with coupling $-1$ gives the principal term $1/8$, and the remaining kernel
vectors are those with parity opposite to \eqref{eq:critical-J}.  This proves
\eqref{eq:threshold-master} for the second determinant.

Finally, with $r=-\delta=m+\frac12$, equation \eqref{eq:threshold-q} gives
$q=(-1)^m(r-m)$.  The algebra in Lemma~\ref{lem:unfolding} remains valid at
this endpoint and turns \eqref{eq:threshold-master} into
\eqref{eq:threshold-unfolded}.
\end{proof}

Proposition~\ref{prop:threshold-asymptotics} completes the proof of
Theorem~\ref{thm:main}.  The argument determines the coefficient of $\log N$;
it deliberately allows additional factors of the form $N^{o(1)}$, including
possible powers of $\log N$.


\section{Example: \texorpdfstring{$\delta=-7/3$}{delta=-7/3}}

For $\delta=-7/3$,
\[
 \alpha=-\frac43,
 \qquad
 s=\sin\left(-\frac{4\pi}{3}\right)=\frac{\sqrt3}{2},
 \qquad
 q=\frac13.
\]
There are two edge eigenvectors: $j=0$ at $+1$ and $j=1$ at $-1$.
For $D_N^-$, the principal-branch contribution is
\[
 -\frac{q^2+q}{2}=-\frac29,
\]
and the $j=0$ tail contributes
\[
 2\delta+1=-\frac{11}{3}.
\]
For $D_N^+$, the principal-branch contribution is
\[
 -\frac{q^2-q}{2}=\frac19,
\]
and the $j=1$ tail contributes
\[
 2\delta+3=-\frac53.
\]
Consequently,
\begin{align*}
 \det(\id_N-H_N^{-7/3})
 &\asymp N^{-35/9},\\
 \det(\id_N+H_N^{-7/3})
 &\asymp N^{-14/9},\\
 \det\bigl(\id_N-(H_N^{-7/3})^2\bigr)
 &\asymp N^{-49/9}.
\end{align*}
Thus the principal Hilbert contribution and the two edge-state tails give not
only the logarithmic exponents but two-sided power bounds.

\section{Exceptional parameters}\label{sec:exceptional}

\subsection{Negative integers}
Let $\delta=-n$, $n\in\mathbb N$.  Although the literal formula
\eqref{eq:def-Hdelta} contains $0/0$ entries, its entrywise analytic
continuation is
\[
 H_N^{-n}(j,k)=
 \begin{cases}
 (-1)^{n-1},&j+k=n-1,\\
 0,&j+k\ne n-1.
 \end{cases}
\]
For $N\ge n$, this is a signed reversal on the first $n$ coordinates and zero
on the remaining coordinates.  Hence $(H_N^{-n})^2$ is the identity on a
nonzero subspace and
\[
 \det\bigl(\id_N-(H_N^{-n})^2\bigr)=0.
\]
The negative integers therefore do not belong to the power-law regime in
Theorem~\ref{thm:main}.

\section{Conclusion}

The main theorem gives two-sided power bounds at every noncritical
parameter and the same powers up to subpolynomial factors at the half-integer
thresholds.  On the left of $\delta=\frac12$ the effective phase is the
unfurled value $-\delta$; on the right it is the principal triangular wave
$q_\delta$.  Thus the squared-determinant exponent is $-\delta^2$ in the
unfolded range and becomes
$-\operatorname{dist}(\delta,\mathbb Z)^2$ for larger positive parameters.

The Gram representation \eqref{eq:gram-representation} shows that the full
parameter $\delta$ is present from the beginning as a displacement of
trigonometric frequencies.  The principal branch appears only after the
shifted generalized Hilbert matrix is compared with the standard Hilbert
matrix, whose determinant asymptotics depend on the scalar coupling
$s_\delta$.

Away from the threshold couplings, that comparison is trace-class stable only
after the exact kernels at $\pm1$ have been regularized.  Removing the
regularization creates a finite-dimensional determinant.
Lemma~\ref{lem:edge-extraction} identifies it, up to bounded factors, with the
Gram determinant of the omitted tails of the edge eigenvectors.  At the
threshold half-integers, the same reduction is made after compressing to the
base parameters $\alpha=\pm\frac12$; the missing uniform gap costs only a
subpolynomial factor.  In both cases the power-law tails contain precisely the
branch index lost on passing from $-\pi\delta$ to
$\arcsin(\sin(-\pi\delta))$.

The resulting relation may be summarized as
\[
 \begin{aligned}
 \text{lifted phase}
 &=\text{principal continuous phase}\\
 &\quad+\text{edge-eigenvector branch data}.
 \end{aligned}
\]
It is analogous to a Levinson-type principle, in which a continuous phase is
supplemented by a bound-state count.  Here both terms are explicit, and the
identity
\[
 -\delta
 =m+(-1)^m\frac1\pi\arcsin(\sin(-\pi\delta))
\]
is recovered directly in the exponents of the finite determinants.

\begin{thebibliography}{99}

\bibitem{AlemanMontesSarafoleanu}
A. Aleman, A. Montes-Rodr\'iguez, and A. Sarafoleanu,
\emph{The eigenfunctions of the Hilbert matrix},
Constr. Approx. \textbf{36} (2012), no.~3, 353--374.
\href{https://doi.org/10.1007/s00365-012-9157-z}{doi:10.1007/s00365-012-9157-z}.

\bibitem{FedeleGebert}
E. Fedele and M. Gebert,
\emph{On determinants identity minus Hankel matrix},
Bull. Lond. Math. Soc. \textbf{51} (2019), no.~4, 751--764.
\href{https://doi.org/10.1112/blms.12271}{doi:10.1112/blms.12271}.

\bibitem{Otte}
P. Otte,
\emph{A Szeg\H{o} limit theorem related to the Hilbert matrix},
Rocky Mountain J. Math. \textbf{54} (2024), no.~5, 1447--1472.
\href{https://doi.org/10.1216/rmj.2024.54.1447}{doi:10.1216/rmj.2024.54.1447}.

\bibitem{RosenblumII}
M. Rosenblum,
\emph{On the Hilbert matrix. II},
Proc. Amer. Math. Soc. \textbf{9} (1958), 581--585.

\bibitem{Silbermann}
B. Silbermann,
\emph{On the spectrum of Hilbert matrix operator},
Integral Equations Operator Theory \textbf{93} (2021), Paper No.~21.
\href{https://doi.org/10.1007/s00020-021-02637-5}{doi:10.1007/s00020-021-02637-5}.

\bibitem{Wilf}
H. S. Wilf,
\emph{Finite sections of some classical inequalities},
Ergebnisse der Mathematik und ihrer Grenzgebiete, Band~52,
Springer-Verlag, Berlin--New York, 1970.
\href{https://doi.org/10.1007/978-3-642-86712-5}{doi:10.1007/978-3-642-86712-5}.

\end{thebibliography}

\end{document}
